\section{Q1 事实核查：Nielsen 说的代码事实是否属实}
\label{sec:facts}

本节所有源码摘录均由脚本从\emph{作者原始源码包}（哈希已核对）按行号截取，未经手工改写；过长的行在报告中截断并标注。

\subsection{C1：恒真的几何谓词——属实}
\label{sec:fact-def}

\excerpt{anygeo}

同一个文件里并排放着两个同类型（\lean{GeoP}）的谓词：\lean{geoOf} 要求存在正整数 $c$ 使 \lean{IsGeoOf} 成立，并且 \lean{HasTrueSurf} 成立；
\lean{anyGeometry} 对任何输入都返回 \lean{True}。全库共有 24 个文件、118 行出现 \lean{anyGeometry}（与外部报道一致）。

谓词 $\Gamma$ 通过下面两个结构进入证明：\lean{OAt} 是带 \lean{geo}、\lean{lam} 两个字段的「证书」，\lean{GeoCl} 是把 $\Gamma$ 沿除子变换封闭起来的归纳谓词。

\excerpt{oat}
\excerpt{geocl}

\subsection{C2：它在主定理的依赖链上吗——在；但点名的 \texttt{plugAll} 不在}
\label{sec:fact-path}

这是公开讨论中悬而未决的关键问题（DEV 文章原话：没有确定 \path{Main.lean} 的依赖路径是否经过 \lean{anyGeometry}）。我们从两个独立的层面回答。

\paragraph{（a）源码层面。}主定理文件只有 12 行：

\excerpt{main}

沿着 \lean{rhOfAll} $\to$ \lean{allChi} $\to$ \lean{UpperTheory.FinalTheorems.everyCharacter} $\to$ \lean{…every_character_of_literature_of_vWin}
一路向下，到达下面这个定理；它的证明三次把 \texttt{geoCl anyGeometry} 作为显式参数传给下游引理：

\excerpt{siegel_a}
\excerpt{siegel_b}

所以 \lean{anyGeometry} 确实在源码层面的调用链上。但请注意 Nielsen 点名的定理 \lean{plugAll}：

\excerpt{plugall}
\excerpt{everyE2}

全库检索表明，\lean{plugAll} 只有一个调用者 \leanu{everyE₂}，而 \leanu{everyE₂} \emph{没有任何调用者}。主路径实际使用的是 \path{ThirdPlaceSeries/} 目录下与之逐句平行的孪生版本
（\lean{HypsAlt} 对应 \lean{Hyps}，\lean{every_character_of_green_v_win} 对应 \lean{plugAll}，\lean{parametersOfExtension} 对应 \lean{prmR}）：

\excerpt{topplug}

两条平行路线同样硬编码了 \texttt{geoCl anyGeometry}，因此这一出入不影响 C2 的实质，只影响其表述的精确性。

\paragraph{（b）内核层面（决定性）。}源码检索只能给出「文本上引用了」的近似结论。为得到精确答案，我们用 conescan 扫描了官方导出文件
（1,080,769,508 行、1,065,214,335 个表达式结点、1,468,418 条声明记录），以 \lean{GlobalTheory.Mem.Top.Bl.anyGeometry} 为种子做污点传播。结果：

\begin{verdictbox}
在五个主定理的依赖锥（1,468,418 个声明）中，\textbf{恰有 30 个声明}传递依赖 \lean{anyGeometry}，五个主定理全部在内。一条内核级见证链是：

{\small\lean{riemannHypothesis} $\to$ \lean{rhOfAll} $\to$ \lean{allChi} $\to$ \lean{UpperTheory.FinalTheorems.everyCharacter} $\to$ \lean{…everyCharacter_of_row}
$\to$ \lean{UpperTheory.LiteratureTheorems.every_character_of_literature_of_vWin} $\to$ \lean{UpperTheory.Lit.every_character_of_theta_integrals_local_matching}
$\to$ \lean{UpperTheory.Lit.nonempty_hyps_definite_side_of_closed_bound} $\to$ \lean{GlobalTheory.Fc.S0.StageU.hypotheses_nonempty_of_unramified_exponents_alt}
$\to$ \lean{GlobalTheory.Mem.Top.Bl.anyGeometry}。}

同一次扫描还表明：\lean{UpperTheory.Top.plugAll}、\lean{plugWin}、\leanu{everyE₂}、\leanu{hypsE₂}、\lean{GlobalTheory.Fc.Hyps}、\lean{prmR}、\lean{Prm}
这些名字\textbf{从未出现在导出文件中}，即它们不在任何主定理的依赖锥内。
\end{verdictbox}

30 个声明的完整清单见附录表~\ref{tab:tainted}。值得注意的是其中「出现在类型里」一列：从 \lean{every_character_of_theta_integrals_local_matching} 往上，
\lean{anyGeometry} 只出现在\emph{证明项}里而不出现在\emph{命题}里——这些定理的陈述已经与 $\Gamma$ 无关。

\subsection{C3：「代码自称 degenerate 的模型」——属实}
\label{sec:fact-degenerate}

DEV 文章称其检索 \texttt{degenerate} 只找到数学术语 nondegenerate，因而未能确认这一条。实际上相关定义采用驼峰命名（大写 D），区分大小写的检索会漏掉它：

\excerpt{geodegen}

这是抽象几何接口 \texttt{DV.Geo F} 的一个「占位」实现：簇取为 $\mathbb{N}$，有效性、nef、big 等谓词一律取 \lean{True}，闭子簇类取单点 \lean{PUnit}。
它所在的文件就叫 \path{Inhabited.lean}，同文件还有 \lean{cxJunk}、\lean{mdlJunk}、\lean{algDegenerate}、\lean{padDegenerate} 等同类定义。
它们经由 \lean{divXFl} 进入主路径上的引理 \lean{geoOver_stages}：

\excerpt{divxfl}
\excerpt{geoover}
\excerpt{geoover2}

第二遍 conescan 以这些名字为种子，确认它们\emph{全部}位于 \lean{riemannHypothesis} 的内核级依赖锥中：

\begin{center}\footnotesize
\begin{tabularx}{\textwidth}{@{}l>{\raggedright\arraybackslash}Xrc@{}}
\toprule
种子 & 导出文件中匹配到的全名 & 依赖它的声明数 & 主定理依赖它？\\
\midrule
\tableinput{tables/scan2_counts.tex}
\bottomrule
\end{tabularx}
\end{center}

{\small 注：\lean{geoDegenerateNoTau} 的名字在导出中出现但没有对应的声明记录、也没有被任何声明引用；\lean{geoTrue}（由「真几何」\lean{geoOf} 构造的 \lean{DV.Geo}）、
\lean{shimUJunk}、\lean{NormalizedHeckeDataPlaceholder} 的名字完全没有出现。末两行是与第一遍扫描的一致性对照。}

\paragraph{与 C1 的一个重要区别。}第三遍扫描表明，抽象几何接口 \lean{Formalism.BasicNotions.DV.Geo} 与 $\Gamma$ 槽的处境并不相同：
依赖锥中有 736 个声明（占 1,468,418 的 0.05\%）以某种方式涉及 \lean{DV.Geo} 类型，其中 621 个（传递地）读取了它的某个字段——
实际被用到的字段投影只有 \lean{Var}、\lean{Dv}、\lean{H0}、\lean{inter}、\lean{vanishOn}、\lean{arVol} 等 18 个。
也就是说，这个接口\emph{不是}无人读取的死槽，而是「一小层对任意几何数据成立的泛型声明，被用在了一个平凡实例上」。
这不改变第~\ref{sec:logic}~节的结论，但两种情形应当分开表述，本报告不把「只写不读」这个更强的说法用在 \lean{geoDegenerate} 上。

\subsection{C4、C5：「按定义成立」与「由 \texttt{rfl} 成立」——机制属实}
\label{sec:fact-rfl}

\excerpt{heighteq}

这段代码先把高度函数\emph{定义}为显式公式 \texttt{ht := DV.hExpl6 …}，再用 \lean{setDiv} 构造一个以它为高度的新数据，于是「高度等于该公式」（\lean{h_eq}）和
「格等于 \lean{nPrime}」（\lean{nprimeDef}）自然由 \lean{rfl} 成立。\lean{LedPart.pt}、\lean{LedPart.ht}、\lean{setDiv} 均在依赖锥内（第二遍扫描确认）。
所以 C5 的描述在机制上是准确的；C4 则只对 \lean{nprimeDef} 这类字段准确——交叉验证代理 C 指出同一见证的其他格性质（如 \lean{NuDef}、\lean{nuLat_comap_isLattice}）是用实质性定理证明的，
这一点我们只做了抽查，未逐项复核。

「先定义对象、再由定义读出它满足的等式」是构造性存在证明的标准写法：要证 $\exists h,\ h = F$，取 $h := F$ 即可。它是否「掏空」了什么，取决于下游是否需要 $h$ 是某个\emph{独立定义}的量——
而这正是第~\ref{sec:logic}~节要回答的问题。

\subsection{C6：「提取几何内容的引理从未被调用」——属实，且可加强为内核级结论}
\label{sec:fact-unused}

源码检索显示：\lean{GeoCl} 这个名字在全库只出现在它的定义文件 \path{NotGeoClFalse.lean} 和一处 \lean{GeoCl.step} 的使用中；
构造子 \lean{GeoCl.base} 只经由引理 \lean{geoCl_of} 被用到一次；全库没有任何写明针对 \lean{GeoCl} 的分情形分析或归纳。
「真几何」谓词 \lean{GlobalTheory.Mem.Top.Bl.geoOf} 只在定义处和 \path{NotLamAreShift.lean} 的三个无人调用的引理中出现
（按单词匹配 \lean{geoOf} 共 10 个文件，其余是不同命名空间下的同名定义）。

但「文本上没有」不等于「证明项里没有」（模式匹配可以不写出类型名）。内核级扫描给出了无歧义的答案：

\begin{verdictbox}
\begin{itemize}
\item \textbf{\lean{GeoCl.rec} 零引用。}在归纳命题上做任何分情形分析或归纳，展开到内核层面都必须经过它的递归子。
导出文件中引用 \lean{GlobalTheory.Mem.Top.Bl.GeoCl.rec} 的只有 \lean{GeoCl} 自己的归纳声明块（4 条记录），辅助定义 \lean{GeoCl.casesOn}、\lean{GeoCl.recOn} 的名字根本没有出现。
也就是说，\emph{依赖锥里没有任何一个证明曾经「打开」过一个 \lean{GeoCl} 证书}。
\item \textbf{\lean{OAt} 的 \lean{geo} 字段只被两个定理读取}（\lean{divEx_geoCl} 与 \lean{oAt_setDiv}），而这两处读取的目的都只是把它原样塞进 \lean{GeoCl} 的构造子里（见下）。
结构 \lean{OAt} 的递归子 \lean{OAt.rec} 同样零引用，所以投影是读取其字段的唯一途径。
作为对照，同一结构的另一个字段 \lean{lam} 除了这两处转手之外，还被 \lean{GlobalTheory.T.parameter_of_discrete_series_ray_alt} 真正使用——
同一个证书里，一个字段有消费者，另一个没有。
\item \textbf{\lean{geoOf}、\lean{IsGeoOf}、\lean{HasTrueSurf} 不在依赖锥内}——这三个名字在导出文件中完全没有出现。
\end{itemize}
\end{verdictbox}

\excerpt{divexprov}
\excerpt{divexgeocl}
\excerpt{oatsetdiv}

\lean{divEx_prov} 是整条路径上真正「兑现」$\Gamma = \mtt{True}$ 的地方：它用 \lean{trivial} 填上 \lean{geo} 字段，同时用实证的 \texttt{lamAre\_heq …} 填上 \lean{lam} 字段。
之后 \lean{geo} 只被转手（\lean{GeoCl.base}、\lean{GeoCl.step}），再也没有被打开。

\subsection{小结}

Nielsen 的事实陈述经得起核查，有的地方（C3、C6）甚至比已有的公开复核所能确认的更扎实。争议的全部重量因此落在下一个问题上：这些事实\emph{意味着}什么。
